\documentclass[11pt]{article}
\usepackage[a4paper,margin=25mm]{geometry}
\usepackage{amsmath,amssymb,amsthm,mathtools,booktabs,hyperref}
\hypersetup{hidelinks}
\emergencystretch=2em
\newcommand{\E}{\mathbb E}
\newcommand{\Prob}{\mathbb P}
\newcommand{\R}{\mathbb R}
\newcommand{\norm}[1]{\lVert #1\rVert}
\newcommand{\dd}{\mathop{\mathrm d}}
\newtheorem{proposition}{Proposition}
\newtheorem{lemma}[proposition]{Lemma}
\newtheorem{corollary}[proposition]{Corollary}
\theoremstyle{remark}
\newtheorem*{remark}{Remark}
\title{Recursive IRMF component bounds: verified partial results}
\author{}
\date{31 October 2026}
\begin{document}
\maketitle

\noindent\textbf{Scope.}
This note checks and develops the valid parts of the submitted contraction
draft. It proves an exact decomposition,
conditional deterministic operator bounds, concentration of the leading term,
and a single-pass remainder bound. The final component bound assumes a
simultaneous propagation-remainder estimate. That estimate, and comparison
with the population targets in Section 3.2 of \texttt{IMF.pdf}, remain open.
The accompanying audit is in
\path{research/recursive-component-bounds/README.md}.

\section{Definitions and target}
Let \(Y=X+\varepsilon\in\R^n\), with deterministic \(X\) and independent,
identically distributed noise coordinates. Take the smoothed absolute loss
\(\rho_H\), with \(H>0\), and set
\[
 |\rho_H'|\le1,\qquad 0<\rho_H''\le L_2,\qquad
 |\rho_H'''|\le L_3.
\]
Assume \(\E|\varepsilon_u|<\infty\), so
\(\phi(c)=\E\rho_H(c+\varepsilon_u)\) is finite. Alternatively, define
\(\phi(c)=\E[\rho_H(c+\varepsilon_u)-\rho_H(\varepsilon_u)]\);
the difference is bounded by \(|c|\), and the same derivatives and minimizers
apply without a moment assumption. All minimizers below are finite and
unique. For the anchored definition this follows from strict convexity and
the score limits at \(\pm\infty\).

For stages \(1\le k\le K\), let \(a^{(k)}_{tu}\ge0\) be deterministic weights
with \(\sum_u a^{(k)}_{tu}=1\), and write
\(\alpha_k=\max_{t,u}a^{(k)}_{tu}\). Define
\[
 \mathcal M_k(y)_t=\arg\min_s\sum_u a^{(k)}_{tu}\rho_H(y_u-s),\qquad
 Y^{(1)}=Y,\quad \widetilde S^{(k)}=\mathcal M_k(Y^{(k)}),\quad
 Y^{(k+1)}=Y^{(k)}-\widetilde S^{(k)}.
\]
The deterministic reference uses the same \(\phi\) at every stage:
\[
 X_\phi^{(1)}=X,\qquad
 S_{\phi,t}^{(k)}=\arg\min_s\sum_u a^{(k)}_{tu}\phi(X_{\phi,u}^{(k)}-s),
 \qquad X_\phi^{(k+1)}=X_\phi^{(k)}-S_\phi^{(k)}.
\]
Set \(e_k=\widetilde S^{(k)}-S_\phi^{(k)}\) and
\(E_k=\sum_{j<k}e_j\). Thus \(Y^{(k)}=X_\phi^{(k)}+\varepsilon-E_k\).
All single passes below use the \emph{same} noise vector \(\varepsilon\).
Independence is across observations, not stages.

With \(c^{(k)}_{tu}=X_{\phi,u}^{(k)}-S_{\phi,t}^{(k)}\), define
\begin{align*}
 f_{kt}&=\sum_u a^{(k)}_{tu}\phi''(c^{(k)}_{tu}),&
 q_{kt}&=\sum_u a^{(k)}_{tu}\rho_H'(c^{(k)}_{tu}+\varepsilon_u),&
 b_{kt}&=q_{kt}/f_{kt},\\
 V_k[t,u]&=a^{(k)}_{tu}\phi''(c^{(k)}_{tu})/f_{kt}.&&
\end{align*}
The population score equation implies \(\E q_{kt}=0\).
The matrices \(V_k\) are deterministic and row stochastic. In the source's
unnormalized notation, \(f_{kt}=F_t^{(k)}/N_k(t)\), \(b_k=F^{(k)-1}\nabla^{(k)}\),
and \(V_k=W_\phi^{(k)}\).

\section{Exact propagation identity}
For any sequence \(V_k\), put
\[
 B_{m:l}=(I-V_m)\cdots(I-V_l),\quad B_{m:l}=I\ (m<l),\qquad
 T_{kk}=I,\quad T_{kj}=-V_kB_{k-1:j+1}\ (j<k).
\]
The order matters because the population kernels need not commute. Define
\begin{align*}
 R_k&=\mathcal M_k(X_\phi^{(k)}+\varepsilon)-S_\phi^{(k)}-b_k,\\
 \Phi_k&=\mathcal M_k(X_\phi^{(k)}+\varepsilon)
            -\mathcal M_k(Y^{(k)})-V_kE_k.
\end{align*}
\begin{proposition}\label{prop:identity}
For every realization,
\[
 e_k=L_k+Q_k,\qquad L_k=\sum_{j\le k}T_{kj}b_j,\qquad
 Q_k=\sum_{j\le k}T_{kj}(R_j-\Phi_j).
\]
\end{proposition}
\begin{proof}
The definitions give \(e_k=b_k+R_k-\Phi_k-V_kE_k\). Hence
\(E_{k+1}=(I-V_k)E_k+b_k+R_k-\Phi_k\), with \(E_1=0\).
Unroll this recurrence and substitute it into the expression for \(e_k\).
\end{proof}
For quadratic loss, \(R_k=\Phi_k=0\) and
\(L_k=W_k(I-W_{k-1})\cdots(I-W_1)\varepsilon\).
For the robust loss, \(\E[\partial b_k/\partial\varepsilon]=V_k\), so
\[
 \E[\partial L_k/\partial\varepsilon]
 =\sum_{j\le k}T_{kj}V_j=V_kB_{k-1:1}.
\]
This is an identity for an expected derivative. In general it does not
represent \(L_k\) as one composite matrix applied to a common score vector.

\section{Deterministic operator estimates}
We use \(\norm{C}_\infty=\max_t\sum_u|C[t,u]|\).
For linear periodic smoothers \(W_k\), let \(w_k(\omega)\) be the real Fourier
multiplier of their convolution kernels on \(\mathbb Z\). Let
\(r_k=nh_k\ge1\), \(r_{k+1}=r_k/a\), \(a>1\). Assume, for fixed \(c,C_-\),
\begin{align}
 |w_k(\omega)|&\le c\min\{1,(r_k|\omega|)^{-2}\},\label{eq:K1}\\
 |w_k'(\omega)|&\le cr_k\min\{r_k|\omega|,(r_k|\omega|)^{-2}\},\label{eq:K2}\\
 0\le1-w_k(\omega)&\le c\min\{(r_k\omega)^2,1\},\label{eq:K3}\\
 \prod_{k=1}^K(1+(w_k(\omega))_-)&\le C_- .\label{eq:K4}
\end{align}
Values at zero are interpreted by continuity. These are hypotheses on the
actual discrete multipliers. A continuous-kernel formula or a finite
frequency plot alone does not verify them uniformly over all designs.

\begin{lemma}\label{lem:linear}
Under \eqref{eq:K1}--\eqref{eq:K4}, there is \(C_*=C(c,C_-,a)\ge1\) such that,
with \(A_{m:l}=(I-W_m)\cdots(I-W_l)\),
\[
 \norm{A_{m:l}}_2\le C_-,\qquad
 \norm{A_{m:l}}_\infty\le C_P:=1+C_*,\qquad
 \norm{W_kA_{k-1:l}}_\infty\le C_* .
\]
Empty products are included.
\end{lemma}
\begin{proof}
Use \(\norm{g}_{L^2}^2=(2\pi)^{-1}\int_{-\pi}^{\pi}|g|^2\dd\omega\).
Weighted Cauchy--Schwarz and Parseval give, for an impulse response \(v\),
\[
 \norm{v}_1\le(1+\pi r)^{1/2}
  \bigl(\norm{\widehat v}_{L^2}^2+r^{-2}\norm{\widehat v'}_{L^2}^2\bigr)^{1/2},
 \qquad r\ge1.
\]
Every subproduct \(\Pi\) of the factors \(1-w_i\) is at most \(C_-\).
For any consecutive index range, the geometric scales imply
\[
 \sum_i|w_i'(\omega)|\le\frac{cG_a}{|\omega|},\qquad
 G_a=\frac1{1-a^{-2}}+\frac1{1-a^{-1}}.
\]
Indeed, \(r_i\min(r_i|\omega|,(r_i|\omega|)^{-2})
=|\omega|^{-1}\min((r_i|\omega|)^2,(r_i|\omega|)^{-1})\), and the
terms below and above one are geometric series.

For \(\widehat v=w_k\prod_{i=l}^{k-1}(1-w_i)\),
\(\norm{\widehat v}_{L^2}^2\le C/r_k\) by \eqref{eq:K1}.
In its derivative, the term containing \(w_k'\) has squared integral
at most \(Cr_k\). The term differentiating \(w_{k-1}\) has squared integral
at most \(Cr_{k-1}=Car_k\). All other differentiated terms retain
\(1-w_{k-1}\). Their sum is bounded by
\[
 C\min\{1,(r_k|\omega|)^{-2}\}
       \min\{a^2r_k^2|\omega|,|\omega|^{-1}\}.
\]
Split the integral at \(1/(ar_k)\) and \(1/r_k\) to bound its squared
integral by \(C(a)r_k\). For an empty product only the \(w_k'\) term remains.
The weighted inequality with \(r=r_k\) proves the component bound.

For \(\widehat g=\prod_{i=l}^m(1-w_i)-1\), telescoping gives
\(|\widehat g|\le C\min\{1,(r_m|\omega|)^{-2}\}\), hence a squared
integral at most \(C/r_m\). In \(\widehat g'\), separate the \(i=m\) term.
All others retain \(1-w_m\); their sum is at most
\(C\min\{r_m^2|\omega|,|\omega|^{-1}\}\).
The squared integral is at most \(Cr_m\). Thus \(\norm g_1\le C_*\) and
\(\norm{I+g}_\infty\le1+C_*\). The spectral norm bound follows from
\eqref{eq:K4}. Periodizing an impulse response cannot increase its
\(\ell_1\) norm, so these bounds also hold on the finite periodic grid.
\end{proof}

\begin{lemma}\label{lem:perturb}
Suppose all linear partial products, including empty ones, obey bounds
\(C_P\ge1\) and \(C_*\ge1\) as above. Set
\(\tau_k=\norm{V_k-W_k}_\infty\), \(\tau_\Sigma=\sum_k\tau_k\).
If \(C_P\tau_\Sigma\le1/2\), then
\[
 \sup_{l\le m}\norm{B_{m:l}}_\infty\le2C_P,\qquad
 \sup_{l\le k}\norm{V_kB_{k-1:l}}_\infty\le2C_*+1.
\]
Moreover \(\tau_k\le\norm{\phi'''}_\infty\omega_k/g_k\), where \(\omega_k\)
is the largest clean-residual oscillation within a stage-\(k\) window and
\(g_k>0\) is the smallest active curvature \(\phi''(c^{(k)}_{tu})\).
\end{lemma}
\begin{proof}
Write \(\Delta_i=V_i-W_i\). The ordered telescoping formula is
\[
 B_{m:l}-A_{m:l}=-\sum_{i=l}^m A_{m:i+1}\Delta_iB_{i-1:l}.
\]
Taking the maximum over the finite collection of partial products yields
\(M\le C_P+C_P\tau_\Sigma M\), so \(M\le2C_P\).
Multiplying the corresponding expansion by \(W_k\), then adding
\(\Delta_kB_{k-1:l}\), bounds the component by
\(C_*+2C_*C_P\tau_\Sigma+2C_P\tau_k\le2C_*+1\).
For one row let \(g_u=\phi''(c^{(k)}_{tu})\), \(\bar g=\sum_u a^{(k)}_{tu}g_u\).
Then \(\sum_u a^{(k)}_{tu}|g_u/\bar g-1|
\le(\max g_u-\min g_u)/\min g_u\le\norm{\phi'''}_\infty\omega_k/g_k\).
\end{proof}

\section{Probability bounds}
Assume \(f_{kt}\ge\lambda>0\),
\(\norm{T_{kj}}_\infty\le C_W\) with \(C_W\ge1\), and
\[
 \alpha_j\le C_\alpha\alpha_k a^{-(k-j)}\quad(j\le k),\qquad C_\alpha\ge1.
\]
The factor \(C_\alpha\) accommodates rounded window schedules. A rate uniform
over a growing number of stages requires \(C_W,C_\alpha,\lambda^{-1}\)
and the remainder constants below to remain bounded. Fix \(x>0\)
and put \(z=x+\log(6nK)\).

\begin{proposition}\label{prop:concentration}
Except on an event of probability at most \(2nKe^{-z}\), simultaneously,
\[
 \norm{L_k}_\infty\le\ell_k:=
 \frac{C_W}{\lambda}\left\{
 \frac{\sqrt{2C_\alpha\alpha_k z}}{1-a^{-1/2}}
 +\frac{2C_\alpha\alpha_k z}{3(1-a^{-1})}\right\}.
\]
\end{proposition}
\begin{proof}
For fixed \(k,t\), collect terms by observation to write
\(L_{kt}=\sum_u\xi_u(\varepsilon_u)\), and define
\[
 d_j(u)=\lambda^{-1}\sum_v|T_{kj}[t,v]|a^{(j)}_{vu}.
\]
Then \(|\xi_u|\le\sum_jd_j(u)\), \(\norm{d_j}_1\le C_W/\lambda\), and
\(\norm{d_j}_\infty\le C_W\alpha_j/\lambda\). Since \(\E L_{kt}=0\),
\(L_{kt}=\sum_u(\xi_u-\E\xi_u)\). These summands are independent,
their total variance is at most
\((C_W/\lambda)^2(\sum_{j\le k}\sqrt{\alpha_j})^2\), and their magnitudes
are at most \(2(C_W/\lambda)\sum_{j\le k}\alpha_j\).
Bernstein's bound \(\sqrt{2vz}+bz/3\) for variance proxy \(v\) and
summand bound \(b\), followed by geometric summation and a union bound,
proves the claim. Individual \(\xi_u\) need not have mean zero.
\end{proof}

\begin{proposition}[Single-pass remainder]\label{prop:fisher}
Put
\[
 u_k=\sqrt{2\alpha_kz}+\tfrac23\alpha_kz,\qquad
 v_k=L_2\bigl(\sqrt{2\alpha_kz}+\tfrac13\alpha_kz\bigr).
\]
Assume \(v_k\le\lambda/4\) and \(u_k\le\lambda^2/(8L_3)\) for every
stage. Except on an event of probability at most \(4nKe^{-z}\),
\[
 \norm{R_k}_\infty\le r_k^{\rm F}:=
 \frac{2u_kv_k}{\lambda^2}+\frac{2L_3u_k^2}{\lambda^3},\qquad k\le K.
\]
If also \(z\alpha_k\le1\), this is at most \(A_Rz\alpha_k\), where
\[
 A_R=\frac{2L_2(\sqrt2+2/3)(\sqrt2+1/3)}{\lambda^2}
       +\frac{2L_3(\sqrt2+2/3)^2}{\lambda^3}.
\]
\end{proposition}
\begin{proof}
For each single pass, Bernstein bounds the score at its population target
by \(u_k\) and the difference between empirical and population curvature
by \(v_k\). Use \(\sum_u(a^{(k)}_{tu})^2\le\alpha_k\), score bound one,
and curvature in \([0,L_2]\), then union over \(nK\) pairs and two statistics.
Throughout distance \(2u_k/\lambda\) from that target, empirical curvature
is at least \(\lambda-v_k-2L_3u_k/\lambda\ge\lambda/2\).
Consequently the empirical score has its zero within this interval.
Write its displacement as \(d\), its score at the target as \(q\), and its
curvature there as \(g\). Taylor's formula gives
\(|q-gd|\le L_3d^2/2\), \(|d|\le2u_k/\lambda\), and \(|g-f_{kt}|\le v_k\).
Therefore
\[
 |d-q/f_{kt}|\le\frac{v_k|d|+L_3d^2/2}{\lambda}
 \le\frac{2u_kv_k}{\lambda^2}+\frac{2L_3u_k^2}{\lambda^3}.
\]
The last assertion follows by bounding \(\alpha_kz\) by
\(\sqrt{\alpha_kz}\). No global strong-convexity assumption is used.
\end{proof}

\section{What follows if propagation is controlled}
\begin{corollary}[Conditional component bound]\label{cor:conditional}
Assume the preceding probability and smallness hypotheses, including
\(z\alpha_k\le1\). Suppose an event \(\Omega_\Phi\), with
\(\Prob(\Omega_\Phi)\ge1-\delta_\Phi\), satisfies
\[
 \Phi_1=0,\qquad \norm{\Phi_k}_\infty\le A_\Phi z\alpha_{k-1}\quad(2\le k\le K).
\]
Then, with probability at least \(1-e^{-x}-\delta_\Phi\), for all stages,
\[
 \norm{e_k}_\infty\le\ell_k+
 \frac{C_WC_\alpha}{1-a^{-1}}(A_R+A_\Phi/a)z\alpha_k.
\]
\end{corollary}
\begin{proof}
The two preceding propositions hold jointly with failure probability at
most \(6nKe^{-z}=e^{-x}\). Intersect with \(\Omega_\Phi\), without requiring
independence. Proposition~\ref{prop:identity} bounds \(\norm{Q_k}_\infty\) by
\(C_W\sum_{j\le k}(\norm{R_j}_\infty+\norm{\Phi_j}_\infty)\).
Now \(\sum_{j\le k}\alpha_j\le C_\alpha\alpha_k/(1-a^{-1})\) and
\(\sum_{j<k}\alpha_j\le C_\alpha\alpha_k/[a(1-a^{-1})]\).
\end{proof}

\noindent The source's propagation lemma does not yet provide this event.
Its quadratic term needs a lower empirical curvature along the whole
interpolation segment. Its fluctuation term needs explicit coefficient
bounds for a degenerate, generally non-identically distributed quadratic
statistic, including the diagonal terms. An induction must also track how
its constants depend on previous remainder bounds and show that the
resulting inequalities close uniformly in the number of stages.

For the paper's targets \(S^{(k)}_{\rm pdf}\), the available statement is only
\[
 \norm{\widetilde S^{(k)}-S^{(k)}_{\rm pdf}}_\infty
 \le\norm{e_k}_\infty+
 \norm{S_\phi^{(k)}-S^{(k)}_{\rm pdf}}_\infty.
\]
The second term is zero at stage one. Its claimed second-order rate at
later stages requires a separate proof. Small Monte Carlo differences
do not establish that rate.

\end{document}
